\documentclass[a4paper, 12pt]{book}

% ---- Language Settings -------------------------------------------------------
\usepackage[english]{babel}     % English
% \usepackage[ngerman]{babel}   % German
% \usepackage[brazilian]{babel} % Brazilian Portuguese
% \usepackage[italian]{babel}   % Italian

% ---- Core Layout & Typography ------------------------------------------------
\usepackage[top=2.5cm, bottom=2.5cm, left=3cm, right=2.5cm]{geometry}
                                % For page dimensions and margin settings
\usepackage{fancyhdr}           % For custom headers and footers
\usepackage{microtype}          % For improved typography, character spacing, and hyphenation
\usepackage{amsmath, amssymb, amsthm}
                                % For advanced math environments, theorems, and symbols
\usepackage{graphicx}           % For including external images and graphics
\usepackage{xcolor}             % For defining and using custom colors

% ---- Text Formatting, Units & Lists ------------------------------------------
\usepackage{enumitem}           % For precise control over list spacing and numbering
\usepackage{nicefrac}           % For readable inline fractions
\usepackage{siunitx}            % For standardized scientific measurements and units
\usepackage{tcolorbox}          % For custom text boxes, callouts, and modular layouts

% ---- Vector Graphics -------------------------------------------------
\usepackage{tikz}               % For generic vector graphics and diagrams
\usepackage{fontawesome5}       % For scalable vector icons

% ---- Music Notation & Layout -------------------------------------------------
\usepackage{guitarchordschemes} % For vector guitar fretboard diagrams
\usepackage{songs}              % For typesetting lyrics with chord anchors

% ---- Hyperref & Cross-Referencing --------------------------------------------
\usepackage{hyperref}           % For clickable links and digital bookmarks
\usepackage{cleveref}           % Must be loaded immediately after hyperref

% ---- Hyperref Setup ----------------------------------------------------------
\hypersetup
{
    colorlinks = true,
    linkcolor  = blue!60!black,
    urlcolor   = red!70!black,
}

% ---- Header/Footer Setup -----------------------------------------------------
\setlength{\headheight}{15pt}
\pagestyle{fancy}
\fancyhf{}

% Override the default chapter and section marks to remove prefix words
\renewcommand{\chaptermark}[1]{\markboth{\thechapter.\ #1}{}}
\renewcommand{\sectionmark}[1]{\markright{\thesection.\ #1}}

% Header: Current Chapter on the left, Current Section on the right
\fancyhead[L]{\small\leftmark}
\fancyhead[R]{\small\rightmark}

% Footer: Page number in the center
\fancyfoot[C]{\thepage}

% ==============================================================================
\begin{document}
% ==============================================================================

\title{Einführung in die Graphentheorie}
\author{Beispielprojekt}
\date{\today}
\maketitle
\tableofcontents
\bigskip

\chapter{Graphentheorie}
% ------------------------------------------------------------------------------
\section{Was ist ein Graph?}
% ------------------------------------------------------------------------------

Ein \emph{Graph} ist eine abstrakte mathematische Struktur,
die Beziehungen zwischen Objekten modelliert.
Sie findet Anwendung in Informatik, Logistik, Sozialwissenschaften
und vielen weiteren Gebieten.

Ein \emph{ungerichteter Graph} ist ein Paar $G = (V, E)$, wobei
\begin{itemize}
\item $V$ eine endliche, nichtleere Menge von \emph{Knoten} (Vertices) und
\item $E \subseteq \binom{V}{2}$ eine Menge von \emph{Kanten} (Edges) ist.
\end{itemize}

Sei $V = \{A, B, C, D\}$ und
$E = \{\{A,B\},\{A,C\},\{B,D\},\{C,D\}\}$.
Der entstehende Graph hat 4 Knoten und 4 Kanten.
Abbildung~\ref{fig:graph} zeigt diese Struktur.

% ---- TikZ-Figur einbinden ----------------------------------------------------
\begin{figure}[ht]
  \centering
  \includegraphics[width=0.45\textwidth]{../figures/build/graph}
  \caption{Einfacher ungerichteter Graph mit 4 Knoten.}
  \label{fig:graph}
\end{figure}

% ------------------------------------------------------------------------------
\section{Grundlegende Eigenschaften}
% ------------------------------------------------------------------------------

Der \emph{Grad} $\deg(v)$ eines Knotens $v \in V$ ist die Anzahl
der Kanten, die $v$ als Endpunkt haben:
$$
\deg(v) = |\{e \in E \mid v \in e\}|.
$$

Für jeden Graphen $G = (V, E)$ gilt:
$$
\sum_{v \in V} \deg(v) = 2|E|.
$$

Jede Kante $\{u, v\} \in E$ trägt zum Grad von $u$ und zum Grad
von $v$ je genau einmal bei, also insgesamt $2|E|$-mal zur Gradsumme.

Im obigen Beispiel: $\deg(A)=2,\;\deg(B)=2,\;\deg(C)=2,\;\deg(D)=2$,
also $\sum \deg = 8 = 2 \cdot 4 = 2|E|$. \checkmark

% ------------------------------------------------------------------------------
\section{Algorithmen auf Graphen}
% ------------------------------------------------------------------------------

Viele klassische Algorithmen operieren auf Graphen.
Abbildung~\ref{fig:bfs} zeigt schematisch den Ablauf einer
\emph{Breitensuche} (BFS).

% ---- TikZ-Figur einbinden ----------------------------------------------------
\begin{figure}[ht]
  \centering
  \includegraphics[width=0.8\textwidth]{../figures/build/bfs}
  \caption{Ablauf der Breitensuche (BFS).}
  \label{fig:bfs}
\end{figure}

\subsection{Breitensuche (BFS)}

Die Zeitkomplexität von BFS ist $\mathcal{O}(|V| + |E|)$,
da jeder Knoten und jede Kante höchstens einmal besucht wird.

% ------------------------------------------------------------------------------
\section{Zusammenfassung}
% ------------------------------------------------------------------------------

Graphen sind ein mächtiges Werkzeug zur Modellierung vernetzter
Strukturen. Das Handshaking-Lemma ist eine fundamentale Eigenschaft,
und Algorithmen wie BFS erlauben eine systematische Traversierung.

Dieser Text dient als Beispiel für den
\textbf{LaTeX Dual-Output Workflow}:
dieselbe \texttt{.tex}-Datei erzeugt sowohl ein druckfertiges PDF
als auch eine HTML-Seite.

\chapter{Numerik}

\section{Einführung}

Numerik ist auch spannend!

% ==============================================================================
\end{document}
% ==============================================================================
